\input{estilo}
% Propuesta, opción 2: la misma entrada que la opción 1, pero lo que no se
% pudo entregar va a una cola aparte, que consume un subconjunto de los
% workers. La principal no se frena: dos colas de distintas velocidades.
% Mismas posiciones que volver-a-encolar.
\begin{document}
\begin{figura}\begin{tikzpicture}
  \newcommand\msj[2]{\node[draw=#1, fill=#1!14, minimum size=4.6mm, inner sep=0pt] at #2 {};}
  \newcommand\flota[4]{% y, ancho, rótulo, nombre
    \foreach \d in {0.3, 0.15} \node[caja, fill=white, minimum width=#2, minimum height=14mm] at (7.2+\d,#1-\d) {};
    \node[caja, minimum width=#2, minimum height=14mm] (#4) at (7.2,#1) {#3};}
  % la principal
  \node[cola, minimum height=40mm, minimum width=11mm] (q) at (2.6,0) {};
  \foreach \i in {1,2,3} {\msj{segundo}{({1.3+0.55*(\i-1)},0)}}
  \draw[msg] (-0.4,0) -- (q.west);
  \flota{0}{22mm}{notifiers}{wp}
  \draw[msg] (q.east) -- (wp.west);
  \node[anchor=west, font=\small] at (0.75,0.9) {principal};
  % la de retry
  \node[cola, minimum height=40mm, minimum width=11mm] (r) at (2.6,-2.4) {};
  \foreach \i in {1,2,3} {\msj{acento}{({1.3+0.55*(\i-1)},-2.4)}}
  \flota{-2.4}{16mm}{algunos}{wr}
  \draw[msg] (r.east) -- (wr.west);
  \node[anchor=west, font=\small] at (0.75,-1.5) {retry};
  \draw[msg, draw=acento] (wp.south west) ++(0.3,0) -- node[etiq, text=acento, left=3pt, pos=0.45] {lo que no pudo entregar} (r.north east);
  \node[etiq, text=suave] at (3.9,-3.55) {la principal no se frena: dos colas, dos velocidades};
  \node[anchor=west] at (-0.4,-4.3) {\textbf{opción 2:} cola aparte de retry};
  \node[anchor=west, font=\footnotesize] at (-0.4,-4.9) {%
    \tikz[baseline=-0.6ex]\node[draw=segundo, fill=segundo!14, minimum size=3.5mm] {};\ su suscriptor responde \quad
    \tikz[baseline=-0.6ex]\node[draw=acento, fill=acento!14, minimum size=3.5mm] {};\ pendiente};
\end{tikzpicture}\end{figura}
\end{document}
